\section{Weil's explicit formula and the heat trace}\label{wt:sec:explicit}

Write the nontrivial zeros as $\rho=\frac12+i\gamma$, with the indexing of the introduction. For an even test function $h$, holomorphic in a closed strip wider than $|\Im r|\le\frac12$ and decaying like $O((1+|\Re r|)^{-2-\eta})$ there, with $g(u)=\frac1{2\pi}\int h(r)e^{-iru}\,dr$, the Guinand and Weil explicit formula \cite{Guinand48,Weil52}, \cite[Theorem 5.12]{IwaniecKowalski04} is
\begin{multline}\label{wt:eq:ef}
\sum_\rho h(\gamma_\rho)=h(\tfrac i2)+h(-\tfrac i2)-g(0)\log\pi\\+\frac1{2\pi}\int_{\R}h(r)\,\Re\psi\big(\tfrac14+\tfrac{ir}2\big)\,dr-2\sum_{n\ge1}\frac{\Lambda(n)}{\sqrt n}\,g(\log n),
\end{multline}
the sum over all nontrivial zeros with multiplicity. The terms $h(\pm\frac i2)$ come from the poles of $\pi^{-s/2}\Gamma(s/2)\zeta(s)$ at $s=1$ and $s=0$, which the factor $s(s-1)$ of $\xi$ removes. Weil's criterion \cite{Weil52,Bombieri} states that RH holds if and only if the right side is nonnegative for every even part of a test function $F(r)\overline{F(\bar r)}$ in the admissible class. Under RH this is clear, since then $h(\gamma)=|F(\gamma)|^2$ for real $\gamma$. The content is the converse. For test functions of small support the positivity is known: Yoshida proved it for support in a short interval \cite{Yoshida92}, and Connes and Consani showed that the archimedean place controls the sign for small support \cite{ConnesConsani21}. % UNCHECKED: Yoshida92 and ConnesConsani21, the statements are the source's and inherit the CHECK status of ch/yo1.tex and ch/dc2.tex.
These are the results of Chapters~\ref{chap:yoshida} and~\ref{chap:zhu}, and they use no information on the zeros.

\subsection*{Gaussian test functions}
Take $h_t(r)=e^{-tr^2}$ with $t>0$. It is entire and decays faster than any power on every strip, so \eqref{wt:eq:ef} holds for every $t>0$, with $g_t(u)=(4\pi t)^{-1/2}e^{-u^2/4t}$. Pairing $\rho$ with $1-\rho$ (the zeros with $\tau>0$ are closed under $\rho\mapsto1-\bar\rho$) gives $\sum_\rho h_t(\gamma_\rho)=2W(t)$ with
\[
W(t)=\sum_je^{-\lambda_jt},\qquad \Re\lambda_j=\tau_j^2-(\tfrac12-\beta_j)^2\ge\tau_j^2-\tfrac14 .
\]
This is the book's unconditional definition of the heat trace (Section~\ref{ch:clock}); the zero counting bound $N(T)=O(T\log T)$ gives absolute convergence for every $t>0$ and after any number of differentiations, and $W$ is real because the nonreal $\gamma_j$ come with their conjugates. Then \eqref{wt:eq:ef} becomes $W(t)=\mathcal A(t)-\mathcal P(t)$, where
\begin{align}
\mathcal A(t)&=e^{t/4}-\tfrac12g_t(0)\log\pi+\frac1{4\pi}\int_{\R}e^{-tr^2}\,\Re\psi\big(\tfrac14+\tfrac{ir}2\big)\,dr,\label{wt:eq:A}\\
\mathcal P(t)&=\sum_{n\ge2}\Lambda(n)\,n^{-1/2}\,g_t(\log n).\label{wt:eq:P}
\end{align}
\begin{resultpar}\label{wt:res:check}
This is the identity \eqref{eq:ch2:Warith} (with the polar term included in $\mathcal A$), and Tables~\ref{tab:ch2:W} and~\ref{tab:ch14:W} give its digits. Table~\ref{wt:tab:check} adds a few more. At 110 digits the value $\mathcal A-\mathcal P$ equals the sum over the first 1700 zeros to $10^{-70}$ or better at each $t$ of the table, $t=0.005,0.02,0.05,0.1,0.2$ (computed, \texttt{code/wt\_heat.py}), and the printed digits of the source are reproduced, for example $W(0.02)=0.018541291731319482254$.\end{resultpar}

\begin{table}[h]\centering\small
\fitwidth{\begin{tabular}{ccc}
\toprule
$t$ & $W(t)$ from the zeros & $\mathcal A(t)-\mathcal P(t)$ from $\Gamma$ and the primes\\
\midrule
0.005 & 0.5371919640910408957 & 0.5371919640910408957\\
0.02 & 0.018541291731319482254 & 0.018541291731319482254\\
0.05 & $4.5878350485735742685\times10^{-5}$ & $4.5878350485735742685\times10^{-5}$\\
0.1 & $2.1047997958223512845\times10^{-9}$ & $2.1047997958223512845\times10^{-9}$\\
0.2 & $4.4301821802236448843\times10^{-18}$ & $4.4301821802236448843\times10^{-18}$\\
\bottomrule
\end{tabular}}
\caption{The explicit formula for the heat trace, agreeing to all printed digits (computed).}\label{wt:tab:check}
\end{table}

\section{Weil positivity is the Thorin condition}\label{wt:sec:thorin}

Put $H(s)=\xi(\frac12+\sqrt s)/\xi(\frac12)$, which is entire in $s$ by the functional equation. By Hadamard's factorisation (order $\frac12$), $H(s)=\prod_j(1+s/\lambda_j)$ with $\sum_j|\lambda_j|^{-1}<\infty$, and the candidate $L=1/H$ has the form
\[
-\log L(s)=\sum_j\log\Big(1+\frac s{\lambda_j}\Big)=\int_0^\infty(1-e^{-sx})\,\frac{W(x)}x\,dx,
\]
by Frullani's integral $\log(1+s/\lambda)=\int_0^\infty(1-e^{-sx})e^{-\lambda x}\,dx/x$ for $\Re\lambda>0$ (the unconditional statement is \cite[Theorem 4.3]{P26}, quoted in Section~\ref{ch:clock}). Under RH every $\lambda_j=\tau_j^2$ is positive, $L(s)=\E e^{-sT}$ for the clock variable $T=\sum_jE_j/\lambda_j$ of Section~\ref{ch:clock} (the reciprocal $\Xi$ law, with $E_j$ independent standard exponentials), its L\'evy measure is $W(x)\,dx/x$ and its Thorin measure is $U=\sum_j\delta_{\lambda_j}$, with the multiplicities. Without RH only the analytic identities hold, and $W(x)/x$ is the only possible L\'evy density (Theorem~\ref{thm:ch12:idsd} proves $W>0$ and $W'<0$ unconditionally).

\begin{theorem}\label{wt:thm:main}
The following are equivalent.
\begin{enumerate}
\item[(i)] RH.
\item[(ii)] Every $\lambda_j$ lies in $(0,\infty)$, that is, the Thorin measure $U=\sum_j\delta_{\lambda_j}$ is supported on $(0,\infty)$ and $T$ is a generalised gamma convolution.
\item[(iii)] $W$ is completely monotone on $(0,\infty)$: $(-1)^nW^{(n)}(t)\ge0$ for all $n\ge0$ and $t>0$.
\item[(iv)] $\mathcal A-\mathcal P$ is completely monotone on $(0,\infty)$, with $\mathcal A$ and $\mathcal P$ as in \eqref{wt:eq:A} and \eqref{wt:eq:P}.
\end{enumerate}
\end{theorem}

\begin{proof}
(iii) and (iv) are the same statement, since $W=\mathcal A-\mathcal P$ for every $t>0$. Under RH every $\lambda_j=\tau_j^2>0$, which gives (ii), and then $W=\int e^{-\lambda t}U(d\lambda)$ is completely monotone (termwise differentiation is justified by the convergence bound above), which gives (iii). For (iii) implies (i) define $S(s)=\sum_j(s+\lambda_j)^{-1}$. The series converges normally away from its poles, and for real $s>\frac14$ Fubini applies because $\sum_j(s+\Re\lambda_j)^{-1}<\infty$, so $S(s)=\int_0^\infty e^{-sx}W(x)\,dx$. If $W$ is completely monotone, Bernstein's theorem gives $W(x)=\int_{[0,\infty)}e^{-ux}\mu(du)$ with $\mu\ge0$, and Tonelli gives $S(s)=\int(s+u)^{-1}\mu(du)$ for $s>\frac14$. Finiteness at one such $s$ gives $\int(1+u)^{-1}\mu(du)<\infty$, so the right side is holomorphic on $\C\setminus(-\infty,0]$. It agrees with the meromorphic series on a real interval, hence on the connected set $\C\setminus(-\infty,0]$ minus the poles of $S$. A pole $-\lambda_j$ outside $(-\infty,0]$ would be removable, but its residue is the positive integer multiplicity of $\lambda_j$. So every $\lambda_j\in[0,\infty)$. Finally $\Im\lambda_j=2\tau_j(\frac12-\beta_j)$ and $\tau_j>0$, so $\beta_j=\frac12$ and $\lambda_j=\tau_j^2>0$. This gives (i), and (ii) implies (iii) as above. (ii) implies (i) by the same last step.
\end{proof}

This is Proposition~\ref{prop:rigid} with the converse spelled out; the uniqueness of exponents of a complex exponential series, which the source asserted, is replaced by the pole argument, which was checked independently (the converse by the resolvent). Statement (iv) is Weil positivity restricted to the cone generated by the Gaussian test functions and their $t$ derivatives, $h(r)=r^{2n}e^{-tr^2}$ for all $n\ge0$ and $t>0$ (these are $|r^ne^{-tr^2/2}|^2$ on the real line and equal $(-\partial_t)^n$ of $h_t$). Positivity of $W$ alone is not this statement, and it is proved without RH (Theorem~\ref{thm:ch12:idsd}). The prime side converges for every $t>0$, because $g_t(\log n)$ decays like $e^{-(\log n)^2/4t}$.

\begin{remark}[The Stieltjes form and the parabola]
The Laplace transform of $W$ is the meromorphic function $S(s)=\sum_j(\lambda_j+s)^{-1}=H'(s)/H(s)$, a Stieltjes function if and only if RH holds. Its Weil kernel $h_s(r)=(r^2+s)^{-1}=|r+i\sqrt s|^{-2}$ is admissible only when the pole $-i\sqrt s$ lies outside the strip, that is $\Re\sqrt s>\frac12$, which is also where the prime sum $\sum\Lambda(n)n^{-1/2-\sqrt s}$ converges. In the $s$ plane this is the outside of the parabola $\Re\sqrt s=\frac12$. A pole $-\gamma^2$ of $S$ with $\gamma=\gamma_r+i\beta$ lies inside the parabola, because $\Re(-\gamma^2)-\frac14+(\Im\gamma^2)^2=-(1-4\beta^2)(\gamma_r^2+\frac14)<0$ for $|\beta|<\frac12$, and for $\beta=0$ the quantity is $-(\gamma_r^2+\frac14)<0$ as well. So all the poles of the Stieltjes picture, violations and non violations alike, sit where the Euler product diverges. The Gaussian form of Theorem~\ref{wt:thm:main} avoids this by using test functions that are admissible everywhere.
\end{remark}

\section{The archimedean place is a Gumbel L\'evy measure}\label{wt:sec:gumbel}

\begin{resultpar}\label{wt:res:acomp}
Put $k(u)=e^{-u/2}/(1-e^{-2u})=\sum_{m\ge0}e^{-(2m+\frac12)u}$. The integral \eqref{wt:eq:A} of the digamma function against the Gaussian can be written in an absolutely convergent form. Start from $\psi(z)=-\gamma_E+\int_0^\infty(e^{-s}-e^{-zs})(1-e^{-s})^{-1}\,ds$ with $z=\frac14+\frac{ir}2$ and $s=2u$, which gives
\begin{equation}\label{wt:eq:Alevy}
\frac1{4\pi}\int e^{-tr^2}\Re\psi\big(\tfrac14+\tfrac{ir}2\big)dr=-\frac{\gamma_E}2g_t(0)+\int_0^\infty\frac{g_t(0)e^{-2u}-e^{-u/2}g_t(u)}{1-e^{-2u}}\,du .
\end{equation}
Splitting off $g_t(0)e^{-u/2}$ from the numerator and using $\int_0^\infty(e^{-u/2}-e^{-2u})(1-e^{-2u})^{-1}du=\frac12(\psi(1)-\psi(\frac14))$, the archimedean part is
\begin{equation}\label{wt:eq:Acomp}
\mathcal A(t)=e^{t/4}+\tfrac12\big[\psi(\tfrac14)-\log\pi\big]g_t(0)+\int_0^\infty\big[g_t(0)-g_t(u)\big]k(u)\,du .
\end{equation}
The compensation $g_t(0)-g_t(u)=O(u^2)$ is essential. The two pieces $\int g_t(u)k(u)\,du$ and $\int g_t(0)k(u)\,du$ separately diverge, because $k(u)\sim1/2u$ at the origin, so the archimedean place cannot be set beside the primes as two finite integrals against $g_t$. The identity \eqref{wt:eq:Alevy} was verified numerically to $10^{-71}$ in the source (unchecked here). The prime part is $\mathcal P(t)=\int_0^\infty e^{-u/2}g_t(u)\,d\psi_{\mathrm C}(e^u)$ with the atom at $u=0$ left out, and both measures have an infinitely divisible meaning.
\begin{itemize}
\item The density $k(u)$ is $u$ times the L\'evy density $e^{-u/2}/\big(u(1-e^{-2u})\big)$ of $-\frac12\log G_{1/4}$, where $G_{1/4}$ has the Gamma$(\frac14)$ law. More generally $-\frac12\log G_a$ has L\'evy density $e^{-2au}/\big(u(1-e^{-2u})\big)$ on $u>0$, with compensation (the density behaves like $1/2u^2$ at the origin), by taking two derivatives of $\log\Gamma(a-i\theta/2)-\log\Gamma(a)$. The variable $-\log G_k$ is the $k$th largest point of a Poisson process with intensity $e^{-v}\,dv$ for an integer $k$ (put $z=e^{-v}$ to get the arrival times of a unit rate process); for a fractional order $a$ it is the log gamma extension of that family, the Gumbel kernel of fractional order of Section~\ref{ev:sec:decomp}, not an ordinary ranked point. It is a real valued infinitely divisible variable, not a subordinator.
\item The prime measure $\Lambda(n)n^{-1/2}\delta_{\log n}$ is $u$ times the atomic measure $\sum_n\frac{\Lambda(n)}{\log n}n^{-\sigma}\delta_{\log n}$ at $\sigma=\frac12$. For $\sigma>1$ this is the L\'evy measure $\sum_{p,m\ge1}p^{-m\sigma}\delta_{m\log p}/m$ of $\log N_\sigma$, with $\mathbb P(N_\sigma=n)=n^{-\sigma}/\zeta(\sigma)$, which is compound Poisson with total mass $\log\zeta(\sigma)$. At $\sigma=\frac12$ the total mass beyond any point is infinite ($\sum_pp^{-1/2}=\infty$), so the measure is locally finite and arithmetic but it is not a L\'evy measure and the law of $N_\sigma$ does not exist. The Gaussian test function is what makes the pairing $\int u\,g_t(u)\,\nu_{1/2}(du)=\mathcal P(t)$ finite.
\end{itemize}\end{resultpar}
So Weil positivity in Gaussian form compares the compensated archimedean functional \eqref{wt:eq:Acomp} with the prime pairing and asks that the difference be completely monotone in $t$. It is not a difference of two probability L\'evy measures.

\subsection*{Fr\'echet form of the prime part}
\begin{resultpar}\label{wt:res:frechet}
As a function of $t$, each prime power enters through
\[
\sqrt{4\pi t}\,g_t(\log n)=\exp\!\Big(-\frac{(\log n)^2}{4t}\Big),
\]
the Fr\'echet distribution function of index one with scale $(\log n)^2/4$, evaluated at $4t/(\log n)^2$ through $\Phi_1(y)=e^{-1/y}$. So
\[
\mathcal P(t)=\frac1{\sqrt{4\pi t}}\sum_{n\ge2}\frac{\Lambda(n)}{\sqrt n}\,\Phi_1\!\Big(\frac{4t}{(\log n)^2}\Big),
\]
a positive weighted series of Fr\'echet distribution functions whose scales are the squared logarithms of the prime powers. The weights $\Lambda(n)n^{-1/2}$ are not summable ($\sum_{n\le N}\Lambda(n)n^{-1/2}\sim2\sqrt N$), so this is not a probability mixture, and a finite truncation is one only after division by its total weight. The smallest scale, $(\log2)^2/4=0.120$, belongs to the prime 2, and for $t\downarrow0$
\[
\mathcal P(t)=\frac{\log2}{\sqrt{8\pi t}}\exp\Big[-\frac{(\log2)^2}{4t}\Big]\Big\{1+O\Big(\exp\Big[-\frac{(\log3)^2-(\log2)^2}{4t}\Big]\Big)\Big\},
\]
proved by isolating $n=2$ and factoring out the exponential of the next term $n=3$. The prime part is governed by this one extreme atom as $t\downarrow0$.\end{resultpar}

\section{Three regimes}\label{wt:sec:regimes}

Table~\ref{wt:tab:regimes} shows $\mathcal A$, $\mathcal P$ and $W$ computed at 110 digits (computed, \texttt{code/\allowbreak wt\_heat.py}; a second run at 130 digits differs by at most $2\times10^{-110}$ in $W$ at all 23 values of $t$ tried). The columns of $\mathcal A$ and $W$ are those of Table~\ref{tab:ch14:W}, where $\mathcal A$ is the sum of its columns $\Pi$ and $\mathcal A$.

\begin{table}[h]\centering\small
\fitwidth{\begin{tabular}{cccccc}
\toprule
$t$ & $\mathcal A(t)$ & $\mathcal P(t)$ & $W(t)$ & digits cancelled & share of the prime 2 in $\mathcal P$\\
\midrule
0.002 & 1.783 & $2.6\times10^{-26}$ & 1.783 & 0 & 1.0000\\
0.005 & 0.537 & $7.2\times10^{-11}$ & 0.537 & 0 & 1.0000\\
0.01 & 0.1497 & $8.4\times10^{-6}$ & 0.1497 & 0 & 1.0000\\
0.02 & 0.02095 & 0.00241 & 0.01854 & 0.05 & 0.9999\\
0.05 & 0.05796 & 0.05791 & $4.59\times10^{-5}$ & 3.1 & 0.966\\
0.1 & 0.16281 & 0.16281 & $2.10\times10^{-9}$ & 7.9 & 0.808\\
0.2 & 0.30152 & 0.30152 & $4.43\times10^{-18}$ & 16.8 & 0.562\\
0.3 & 0.39603 & 0.39603 & $9.32\times10^{-27}$ & 25.6 & 0.427\\
\bottomrule
\end{tabular}}
\caption{Archimedean part, prime part and heat trace (computed). Digits cancelled is $\log_{10}(\mathcal A/W)$, which differs from the measure of the caption of Table~\ref{tab:ch14:W}.}\label{wt:tab:regimes}
\end{table}

\begin{enumerate}
\item Archimedean dominance, $t\le0.01$. Here $\mathcal P/W\le6\times10^{-5}$ (computed: $5.6\times10^{-5}$ at $t=0.01$), and all of $\mathcal P$ comes from the prime 2, through its Fr\'echet factor $e^{-(\log2)^2/4t}$. Positivity of $W$ is inherited from the archimedean part, and so are the derivative conditions of the orders up to 6 tested at the grid of $t$ (unchecked as a computation). The book proves more than this: every order $k\ge1$ for $t\le0.019$ with no information on the zeros (Theorem~\ref{thm:ch12:cmint}(a)), and all orders up to $1.16\times10^{25}$ for every $t>0$ (Theorem~\ref{thm:ch12:window}). A small $\mathcal P/W$ alone controls function values, not derivatives, because the derivatives of $\mathcal P$ are not dominated by those of $\mathcal A$ uniformly in the order. This regime is the Gaussian counterpart of small support positivity: a Gaussian of variance $2t$ sees essentially nothing beyond $\log2$.
\item Crossover, $t\approx0.02$ to $0.05$. The prime part becomes comparable to $W$ (at $t=0.05$ the ratio $\mathcal P/W$ is about $1260$) but is still almost entirely the prime 2. Positivity here compares one extreme atom with the archimedean density.
\item Cancellation, $t\gtrsim0.05$. As $t\to\infty$ both $\mathcal A$ and $\mathcal P$ grow like $e^{t/4}$; on the range of the table $\mathcal A$ rises only from $0.02$ to $0.40$. While $W\approx e^{-\gamma_1^2t}$ (for large $t$, if $\lambda_1$ is simple and separated), the two parts agree to $\log_{10}(\mathcal A/W)\approx\gamma_1^2t/\ln10+\log_{10}\mathcal A\approx87t$ digits ($\gamma_1^2/\ln10=86.767$), and the prime 2's share falls as the bulk of the primes enters. All the information about the zeros lives in this exponentially small difference, and the evaluation of $W$ is very ill conditioned there; this is a statement about conditioning, not a theorem that the information is confined to one interval of $t$.
\end{enumerate}

\subsection{The archimedean positivity window}\label{wt:sec:archwindow}

How far does the archimedean part alone satisfy the Bernstein condition? Combining the two archimedean terms of \eqref{wt:eq:A},
\begin{equation}\label{wt:eq:Aspec}
\mathcal A(t)=e^{t/4}+\int_0^\infty e^{-\lambda t}w(\lambda)\,d\lambda,\qquad w(\lambda)=\frac{\Re\psi(\frac14+\frac i2\sqrt\lambda)-\log\pi}{4\pi\sqrt\lambda}
\end{equation}
(substitute $\lambda=r^2$). So $\mathcal A$ is the Laplace transform of a signed measure: an atom of weight one at $\lambda=-\frac14$, coming from the poles of $\pi^{-s/2}\Gamma(s/2)\zeta(s)$ at $s=1$ and $s=0$, and a density $w$ that is negative for $\lambda<\lambda_*=r_*^2$ and positive above (Lemma~\ref{wt:lem:w}). Here $r_*$ solves $\Re\psi(\frac14+\frac{ir}2)=\log\pi$. Since $\Re\psi(\frac14+\frac{ir}2)-\log\pi=2\vartheta'(r)$ for the Riemann Siegel theta function $\vartheta$ (Section~\ref{ch:heat}), $r_*$ is the critical point of $\vartheta$ (its minimum on $r>0$). Numerically $r_*=6.28984$ and $\lambda_*=39.5620$ (computed, $r_*=6.2898359888$; $r_*\approx2\pi$ because $\Re\psi(\frac14+\frac{ir}2)\approx\log(r/2)$). Only the digamma piece is a Gumbel L\'evy density; the atom and the term in $\log\pi$ are not.

\begin{resultpar}\label{wt:res:window}
The $n$th derivative weights $w$ by $\lambda^ne^{-\lambda t}$, which peaks at $\lambda=n/t$. So $(-1)^n\mathcal A^{(n)}(t)\ge0$ should hold while $n/t$ stays above $\lambda_*$, that is for $t\lesssim n/\lambda_*$ (heuristic, made precise in Proposition~\ref{wt:prop:window}). Table~\ref{wt:tab:window} gives the first sign change $t_n$ of $(-1)^n\mathcal A^{(n)}$, found by bisection at 35 and 45 digits (computed, \texttt{code/wt\_window.py}; the two runs agree). For $n=1,\dots,6$ it is $0.025332,0.051040,0.076135,0.101334,0.126561,0.151805$, within $1$ percent of $n/\lambda_*$ with $\lambda_*=39.5620$.\end{resultpar}

\begin{table}[h]\centering\small
\fitwidth{\begin{tabular}{cccc}
\toprule
order $n$ & first sign change $t_n$ of $(-1)^n\mathcal A^{(n)}$ & $n/\lambda_*$ & relative deviation\\
\midrule
1 & 0.025332 & 0.025277 & 0.22\%\\
2 & 0.051040 & 0.050554 & 0.96\%\\
3 & 0.076135 & 0.075830 & 0.40\%\\
4 & 0.101334 & 0.101107 & 0.22\%\\
5 & 0.126561 & 0.126384 & 0.14\%\\
6 & 0.151805 & 0.151661 & 0.095\%\\
\bottomrule
\end{tabular}}
\caption{The first sign change $t_n$ of $(-1)^n\mathcal A^{(n)}$ against $n/\lambda_*$, $\lambda_*=r_*^2=39.5620$ (computed).}\label{wt:tab:window}
\end{table}

The law $t_n\approx n/\lambda_*$ holds to within 1\% for $n\le6$. The relative deviation is $0.22\%$ at $n=1$ and $0.96\%$ at $n=2$ and decreases from $n=2$ on; in absolute terms $t_n-n/\lambda_*$ is $0.55,4.86,3.05,2.27,1.77,1.44$ in units of $10^{-4}$. The table tests the single inequality $k=n$. Since $t_1<t_2<\dots<t_6$, the conditions of all orders $k\le n$ hold together only on $(0,t_1]$, so $\mathcal A$ is completely monotone of order $n$ only up to $t_1$ and not up to $t_n$. For the full heat trace all six orders are positive at every $t$ tested from $0.002$ to $0.2$ (unchecked as a computation), as Theorem~\ref{thm:ch12:window} proves for all $t$. The binding order is $n=1$: the archimedean part alone is decreasing only up to $t_1=0.02533\approx1/(4\pi^2)=0.025330$, and beyond that point $W$ decreases only because of the prime part, which at $t_1$ is still essentially the prime 2 alone, through $e^{-(\log2)^2/4t_1}=0.0087$ (computed, $0.0087253$).

This gives a spectral reading under RH of the division of labour. The archimedean place supplies the positive spectral density above $\lambda_*\approx4\pi^2$ and two defects: negative spectral mass below $\lambda_*$ and the atom at $-\frac14$. Subtracting the Laplace representations of $\mathcal A$ and $W=\int e^{-\lambda t}U(d\lambda)$ gives $\mathcal P(t)=\int e^{-\lambda t}\,[\delta_{-1/4}+w(\lambda)\,d\lambda-U(d\lambda)]$, so the prime part is not itself a positive Thorin measure; it cancels the defects and leaves the atoms of $U$ at $\lambda_j$, all of which lie above $\lambda_*$ since $\gamma_1^2=199.79$ (computed). Inside $t<t_n$ the single condition $(-1)^n\mathcal A^{(n)}\ge0$ holds without the primes, which is a Gaussian analogue of the small support theorems.

\section{Zeros from primes by Prony's method}\label{wt:sec:prony}

In the cancellation regime $W(t)=\sum_je^{-\lambda_jt}$ is a sum of decaying exponentials, so Prony's method applies, and the book uses it for exactly this purpose with the digit law below (Section~\ref{sec:ch18:reach}, Tables~\ref{ev:tab:prony} and~\ref{ev:tab:prony2}). The difference here is that $W$ comes from \eqref{wt:eq:A} and \eqref{wt:eq:P}, from the gamma function and the primes, with no evaluation of $\zeta$ and no zeros. We computed $W$ at 110 digits at $t=0.02+0.006j$, $j=0,\dots,15$, and fitted eight exponentials. The rule that terms with $\log n>\sqrt{4t\cdot110\ln10}$ are below the working precision keeps the prime powers up to $2767$ at $t=0.062$ (403 primes, 438 prime powers) but up to $38\,436$ at the largest sample $t=0.11$. Cutting at $n\le2767$ at every sample leaves an absolute error $1.2\times10^{-62}$ at $t=0.11$ (against $W=2.85\times10^{-10}$), and the fitted exponents $k\ge4$ are unchanged by it, while the errors of $\gamma_1,\gamma_2,\gamma_3$ are then $4\times10^{-52}$, $5\times10^{-43}$ and $2\times10^{-34}$ (computed). So the recovery of Table~\ref{wt:tab:primezeros} is valid from the 403 primes, and the 110 digit evaluation of $W$ itself needs the larger cutoff at the larger $t$.

\begin{table}[h]\centering\small
\fitwidth{\begin{tabular}{cccc}
\toprule
$k$ & $\widehat\gamma_k$ from primes & $\gamma_k$ & error\\
\midrule
1 & 14.13472514173469379 & 14.13472514173469379 & $1.3\times10^{-52}$\\
2 & 21.022039638771554993 & 21.022039638771554993 & $1.6\times10^{-41}$\\
3 & 25.010857580145688763 & 25.010857580145688763 & $2.9\times10^{-34}$\\
4 & 30.42487612585951321 & 30.42487612585951321 & $7.2\times10^{-24}$\\
5 & 32.935061587739189691 & 32.935061587739189691 & $3.0\times10^{-19}$\\
6 & 37.58617815884 & 37.58617815883 & $1.5\times10^{-11}$\\
7 & 40.9187208 & 40.9187190 & $1.8\times10^{-6}$\\
8 & 43.3279 & 43.3271 & $7.9\times10^{-4}$\\
\bottomrule
\end{tabular}}
\caption{Zeros recovered by Prony's method from the heat trace computed from $\Gamma$ and the primes (computed, \texttt{code/wt\_heat.py}, 110 digits; the source printed errors of $10^{-30}$ for $k\le3$, which was the floor of its reference, and the digits shown are those of the source).}\label{wt:tab:primezeros}
\end{table}

That the explicit formula carries the zeros in its prime side is classical \cite{Guinand48,Weil52}. The zeros are the exponential rates of $W$, the omitted prime atoms act as a bias on the fit, and the number of digits needed to see the $k$th zero grows like $(\gamma_k^2-\gamma_1^2)t/\ln10$ (heuristic, as in Section~\ref{sec:ch18:reach}); this does not count the cost of evaluating $W$ itself.

\section{The prime point process: first and second order}\label{wt:sec:pp}

On the log scale, $u=\log y$, the $\Lambda$ weighted prime powers form a point process with counting function $\psi_{\mathrm C}(e^u)$. By the prime number theorem its intensity is $e^u$, which is the intensity $e^{-v}$ of the Gumbel point process with $v=-u$. In this heuristic reading the prime side of the explicit formula is, to first order, the reflected Gumbel process, and the weight $e^{-u/2}$ in \eqref{wt:eq:ef} is the square root of that intensity, the normalisation under which a Poisson process with intensity $e^u$ has fluctuations of order one. This is an analogy for the cumulative counts and not an identification of the processes.

RH lives in the second order. By von Koch's theorem \cite{vonKoch01}, RH holds if and only if $\psi_{\mathrm C}(y)=y+O(y^{1/2+\epsilon})$ for every $\epsilon>0$, that is $(\psi_{\mathrm C}(y)-y)/y=O(y^{-1/2+\epsilon})$. The error changes sign infinitely often (it is $\Omega_\pm(\sqrt y\log\log\log y)$ by Littlewood), so this is a growth bound and not an index of regular variation. We compare the primes with the Cram\'er random model, which keeps $n$ independently with probability $1/\log n$.
\begin{itemize}
\item For $10^3<y\le10^7$ the normalised fluctuation $|\psi_{\mathrm C}(y)-y|/\sqrt y$ of the primes is at most $0.777$ at integers (attained at $n=1423$) and $0.803$ as a supremum over real $y$ (computed, \texttt{code/wt\_rig.py}, sieve to $2.1\times10^7$). The book quotes the explicit bound $0.94\sqrt y$ in Section~\ref{bp:sec:heat}.
\item For the Cram\'er model the same maximum over $20$ realisations lies between $5.7$ and $12.6$, with median $7.8$ (computed); the source's single draw $11.1$ is within this range.
\end{itemize}
The primes are far more rigid than a Poisson process with the same intensity at this scale. The size of $\psi_{\mathrm C}(y)-y$ at scale $y$ is set under RH by the sum $\sum1/|\rho|$ over the low zeros and by Cram\'er's mean square theorem (Proposition~\ref{wt:prop:energy}), not by repulsion between pairs of zeros; pair correlation governs the variance in short intervals, treated in Section~\ref{wt:sec:rigidity}. Calling the rigidity the prime side counterpart of the repulsion of the zeros is a heuristic.
